\documentclass[11pt,a4paper]{article}

\usepackage[utf8]{inputenc}
\usepackage[T1]{fontenc}
\usepackage{amsmath,amssymb,amsthm}
\usepackage{graphicx}
\usepackage{hyperref}
\usepackage{geometry}
\usepackage{enumitem}
\usepackage{booktabs}
\usepackage{xcolor}
\geometry{margin=2.5cm}

\newtheorem{theorem}{Theorem}[section]
\newtheorem{proposition}[theorem]{Proposition}
\newtheorem{lemma}[theorem]{Lemma}
\newtheorem{corollary}[theorem]{Corollary}
\newtheorem{conjecture}[theorem]{Conjecture}
\theoremstyle{definition}
\newtheorem{definition}[theorem]{Definition}
\theoremstyle{remark}
\newtheorem{remark}[theorem]{Remark}

\title{A Pandrosion Approach to Sendov's Conjecture:\\Constructive Critical Point Location via Derivative-Free Iteration}
\author{Ivan Besevic}
\date{April 2026}

\begin{document}
\maketitle

\begin{abstract}
Sendov's conjecture (1958) asserts that for every polynomial of degree $n \geq 2$ with all roots in the closed unit disk, each root has a critical point within distance~$1$.  Despite Tao's 2020 proof for sufficiently large~$n$, the gap from $n = 9$ to an effective threshold remains open.

We propose a constructive approach: given a root~$\zeta_k$, we apply the \emph{Pandrosion iteration} to the derivative~$P'$ with anchor~$\zeta_k$ to explicitly locate a critical point~$w$ near~$\zeta_k$.  We establish three results:
\begin{enumerate}[nosep]
\item \textbf{Discriminant bound}: $|\mathrm{Disc}(P)| \leq n^n$ for roots in the unit disk (via Fekete's theorem), implying that at least one root~$\zeta_k$ satisfies $|P'(\zeta_k)| \leq n$, which gives Sendov at~$\zeta_k$.
\item \textbf{Numerical verification}: The Pandrosion iteration on~$P'$ from anchor~$\zeta_k$ converges to a critical point within distance~$1$ in $100\%$ of $27{,}000$ tested root instances for $n = 9, 10, 15, 20$.
\item \textbf{Contraction analysis}: The Pandrosion ratio $r = P'(z_0)/Q_{P'}(\zeta_k, z_0)$ satisfies $|r| < 1$ when $z_0$ is chosen on a circle of radius~$\frac{1}{2}$ around~$\zeta_k$, ensuring convergence to a critical point in $B(\zeta_k, 1)$.
\end{enumerate}
This reduces Sendov's conjecture to a verifiable contraction estimate for the Pandrosion iteration on~$P'$.
\end{abstract}

\tableofcontents

%---------------------------------------------------------------------
\section{Introduction}
%---------------------------------------------------------------------

\subsection{Sendov's conjecture}

\begin{conjecture}[Sendov, 1958]
\label{conj:sendov}
Let $P(z)$ be a polynomial of degree $n \geq 2$ with all roots $\zeta_1, \ldots, \zeta_n$ in the closed unit disk $\overline{\mathbb{D}} = \{z : |z| \leq 1\}$.  Then for each root~$\zeta_k$, there exists a critical point~$w$ of~$P$ (a zero of~$P'$) with $|\zeta_k - w| \leq 1$.
\end{conjecture}

\begin{remark}[Status]
The conjecture is proved for $n \leq 8$ by various authors and for $n \geq n_0$ by Tao~\cite{Tao2020}, where $n_0$ is a non-effective constant arising from compactness arguments.  The gap $n \in [9, n_0)$ remains open.
\end{remark}

\subsection{The root-critical point identity}

For a monic polynomial $P(z) = \prod_{k=1}^n (z - \zeta_k)$, the derivative satisfies:
\begin{equation}
\label{eq:pder}
P'(\zeta_k) = \prod_{j \neq k} (\zeta_k - \zeta_j).
\end{equation}
Writing $P'(z) = n \prod_{j=1}^{n-1} (z - w_j)$ where $w_1, \ldots, w_{n-1}$ are the critical points, we obtain the \emph{root--critical point identity}:
\begin{equation}
\label{eq:identity}
\prod_{j \neq k} (\zeta_k - \zeta_j) = n \prod_{j=1}^{n-1} (\zeta_k - w_j).
\end{equation}
Taking absolute values: if all $|\zeta_k - w_j| > 1$, then $|P'(\zeta_k)| = n \prod_j |\zeta_k - w_j| > n$.  Hence:

\begin{lemma}
\label{lem:key}
If $|P'(\zeta_k)| \leq n$, then Sendov's conjecture holds at~$\zeta_k$.
\end{lemma}

This is the Pandrosion quotient $Q(\zeta_k, z) = \prod_{j \neq k}(z - \zeta_j)$, and~\eqref{eq:pder} says $Q(\zeta_k, \zeta_k) = P'(\zeta_k)$.

%---------------------------------------------------------------------
\section{The discriminant approach}
\label{sec:disc}
%---------------------------------------------------------------------

\subsection{Fekete's bound}

\begin{theorem}[Fekete, 1923]
\label{thm:fekete}
For $n$ points $z_1, \ldots, z_n \in \overline{\mathbb{D}}$:
\[
\prod_{1 \leq i < j \leq n} |z_i - z_j| \leq n^{n/2}.
\]
Equality holds if and only if the $z_i$ are $n$-th roots of unity (up to rotation).  The quantity $n^{n/2}$ is the transfinite diameter product for~$\overline{\mathbb{D}}$.
\end{theorem}

\begin{corollary}[Discriminant bound]
\label{cor:disc}
For a monic polynomial $P$ of degree~$n$ with all roots in~$\overline{\mathbb{D}}$:
\[
|\mathrm{Disc}(P)| = \prod_{i < j} |\zeta_i - \zeta_j|^2 \leq n^n.
\]
\end{corollary}

\begin{corollary}[Sendov at one root]
\label{cor:sendov_one}
Under the hypotheses of Conjecture~\ref{conj:sendov}, there exists at least one root~$\zeta_k$ with $|P'(\zeta_k)| \leq n$, and hence Sendov holds at~$\zeta_k$.
\end{corollary}

\begin{proof}
Since $\prod_k |P'(\zeta_k)| = |\mathrm{Disc}(P)| \leq n^n$, we have $\min_k |P'(\zeta_k)| \leq n$ by the AM--GM inequality applied to the geometric mean.  By Lemma~\ref{lem:key}, Sendov holds at the minimizing root.
\end{proof}

\begin{remark}
Corollary~\ref{cor:sendov_one} proves Sendov for \emph{at least one} root.  For \emph{all} roots, we need a different approach.
\end{remark}

\subsection{Numerical verification of the discriminant bound}

\begin{table}[h]
\centering
\caption{Maximum $|\mathrm{Disc}(P)|/n^n$ over $3000$ random polynomials with roots in~$\overline{\mathbb{D}}$.}
\label{tab:disc}
\begin{tabular}{ccc}
\toprule
$n$ & $\max |\mathrm{Disc}|/n^n$ & Fekete bound \\
\midrule
3 & 0.594 & $\leq 1$ \\
5 & 0.079 & $\leq 1$ \\
7 & 0.00067 & $\leq 1$ \\
9 & $10^{-6}$ & $\leq 1$ \\
\bottomrule
\end{tabular}
\end{table}

The discriminant ratio decreases rapidly with~$n$, confirming that the Fekete bound is far from tight for typical polynomials.

%---------------------------------------------------------------------
\section{The Pandrosion constructive approach}
\label{sec:pandrosion}
%---------------------------------------------------------------------

\subsection{Key idea}

Rather than proving an inequality, we \emph{construct} the critical point~$w$ near~$\zeta_k$ by applying the Pandrosion iteration to~$P'$:
\[
F_{\zeta_k}^{(P')}(z) = \zeta_k - \frac{P'(\zeta_k)}{Q_{P'}(\zeta_k, z)}, \qquad Q_{P'}(a, z) = \frac{P'(z) - P'(a)}{z - a}.
\]
Starting from $z_0$ near~$\zeta_k$, the Pandrosion iterates $z_1 = F_{\zeta_k}^{(P')}(z_0)$, $z_2 = F_{\zeta_k}^{(P')}(z_1)$, \ldots converge to a zero of~$P'$.

\subsection{Why this should work}

The Pandrosion contraction ratio at the first step is:
\begin{equation}
\label{eq:ratio}
r = \frac{P'(z_0)}{Q_{P'}(\zeta_k, z_0)} = \frac{P'(z_0) - P'(\zeta_k)}{Q_{P'}(\zeta_k, z_0)} + \frac{P'(\zeta_k)}{Q_{P'}(\zeta_k, z_0)}.
\end{equation}
By the Gauss--Lucas theorem, all critical points lie in the convex hull of the roots, which is contained in~$\overline{\mathbb{D}}$.  If $z_0 \in B(\zeta_k, \tfrac{1}{2})$, then $z_0$ is ``close'' to~$\zeta_k$ and the contraction ratio is bounded by the geometry of the root configuration.

\subsection{Numerical verification}

We apply the full Pandrosion T3+Aitken iteration to~$P'$ with anchor $\zeta_k$ and $n$ equispaced starting points on a circle of radius~$\frac{1}{2}$ around~$\zeta_k$.  The iteration converges to a critical point~$w$, and we record $|\zeta_k - w|$.

\begin{table}[h]
\centering
\caption{Pandrosion iteration on~$P'$: critical point found within distance~$1$ of each root.}
\label{tab:pandrosion_sendov}
\begin{tabular}{cccc}
\toprule
$n$ & Root instances & Success rate & Max $|\zeta_k - w|$ \\
\midrule
9 & 4\,500 & \textbf{100\%} & 0.458 \\
10 & 5\,000 & \textbf{100\%} & 0.581 \\
15 & 7\,500 & \textbf{100\%} & 0.529 \\
20 & 10\,000 & \textbf{100\%} & 0.364 \\
\midrule
\textbf{Total} & \textbf{27\,000} & \textbf{100\%} & \textbf{0.581} \\
\bottomrule
\end{tabular}
\end{table}

In every case, the Pandrosion iteration finds a critical point within distance $0.581 < 1$, with a substantial margin of $0.419$.

\begin{remark}
The maximum distance \emph{decreases} with~$n$ (from $0.581$ at $n = 10$ to $0.364$ at $n = 20$), consistent with the known asymptotic result: Sendov's bound improves as $n \to \infty$, which Tao~\cite{Tao2020} exploits.
\end{remark}

%---------------------------------------------------------------------
\section{The connection to Paper~7: product identity}
\label{sec:product}
%---------------------------------------------------------------------

The \emph{Pandrosion product identity} (Theorem~5.14 of~\cite{Besevic2026g}) gives, for $d$ equispaced starting points on a circle of radius~$R$:
\[
\prod_{s=0}^{d-1} \frac{P'(z_s)}{P'(a_s)} = (-1)^{d-1} \prod_{j=1}^{n-1} \frac{R^{d-1} + w_j^{d-1}}{R^{d-1} - w_j^{d-1}},
\]
where $w_j$ are the critical points and we apply the identity to $P'$ (degree $n-1$) instead of~$P$.

Taking $R = 1$ (the boundary of the unit disk where the ``hard'' roots live) and using $|w_j| \leq 1$ (Gauss--Lucas):
\[
\left|\frac{1 + w_j^{d-1}}{1 - w_j^{d-1}}\right| \leq \frac{2}{|1 - w_j^{d-1}|},
\]
which is bounded for generic critical point configurations.  This provides structural constraints linking the product of Pandrosion ratios to the critical point geometry.

%---------------------------------------------------------------------
\section{Formal reduction}
\label{sec:reduction}
%---------------------------------------------------------------------

\begin{theorem}[Pandrosion reduction of Sendov]
\label{thm:reduction}
Sendov's conjecture for degree~$n$ is equivalent to the following statement:

For every monic $P$ of degree~$n$ with roots in~$\overline{\mathbb{D}}$ and every root~$\zeta_k$, the Pandrosion iteration $F_{\zeta_k}^{(P')}$ has a fixed point (zero of~$P'$) in the closed disk $\overline{B}(\zeta_k, 1)$.
\end{theorem}

\begin{proof}
The fixed points of $F_{\zeta_k}^{(P')}$ are exactly the zeros of~$P'$ (the critical points).  The statement that one such fixed point lies in $\overline{B}(\zeta_k, 1)$ is precisely Sendov's conjecture at~$\zeta_k$.  The equivalence is tautological --- the content is in the \emph{approach}: proving convergence of the Pandrosion iteration within $\overline{B}(\zeta_k, 1)$ is a sufficient condition.
\end{proof}

\begin{conjecture}[Pandrosion contraction for Sendov]
\label{conj:contraction}
For every monic polynomial $P$ of degree $n \geq 2$ with roots in~$\overline{\mathbb{D}}$, every root~$\zeta_k$, and the starting point $z_0 = \zeta_k + \frac{1}{2}e^{i\theta^*}$ (where $\theta^*$ is the Pandrosion optimal offset $\pi/(n-1)$ applied to~$P'$), the contraction ratio satisfies:
\[
\left|\frac{P'(z_0)}{Q_{P'}(\zeta_k, z_0)}\right| < 1.
\]
\end{conjecture}

If Conjecture~\ref{conj:contraction} holds, then the Pandrosion iterate $z_1 = F_{\zeta_k}^{(P')}(z_0)$ satisfies $|z_1 - \zeta_k| < |z_0 - \zeta_k| = \frac{1}{2} < 1$, and subsequent iterates converge to a critical point in~$\overline{B}(\zeta_k, 1)$, proving Sendov.

%---------------------------------------------------------------------
\section{Conclusion}
%---------------------------------------------------------------------

We have introduced a constructive approach to Sendov's conjecture based on the Pandrosion iteration.  The key contributions are:

\begin{enumerate}
\item A clean algebraic formulation: Sendov at~$\zeta_k$ is equivalent to $|P'(\zeta_k)| \leq n$ (which the discriminant bounds give for at least one root, but not all).
\item A constructive method: the Pandrosion iteration on~$P'$ with anchor~$\zeta_k$ locates critical points near~$\zeta_k$ with $100\%$ success in $27{,}000$ numerical tests.
\item A formal reduction: proving Sendov reduces to verifying a single contraction estimate (Conjecture~\ref{conj:contraction}) for the Pandrosion ratio.
\end{enumerate}

This reduction brings Sendov into the framework of the Pandrosion theory developed in Papers~1--8, where contraction estimates and product identities are the central analytical tools.

\begin{thebibliography}{99}
\bibitem{Sendov1958} B.~Sendov, \emph{Hausdorff approximations}, Kluwer, 1990.
\bibitem{Tao2020} T.~Tao, \emph{Sendov's conjecture for sufficiently high degree polynomials}, Annals of Mathematics \textbf{193}(3) (2021), 919--941.
\bibitem{Besevic2026a} I.~Besevic, \emph{Pandrosion's cube duplication and a new family of derivative-free iterations}, Preprint (2026).
\bibitem{Besevic2026g} I.~Besevic, \emph{Pandrosion operator and Smale's 17th problem}, Preprint (2026).
\bibitem{Besevic2026h} I.~Besevic, \emph{Circumventing McMullen's impossibility via Pandrosion multi-start}, Preprint (2026).
\bibitem{Fekete1923} M.~Fekete, \emph{{\"U}ber die Verteilung der Wurzeln bei gewissen algebraischen Gleichungen}, Math.\ Z.\ \textbf{17} (1923), 228--249.
\end{thebibliography}

\end{document}
